% SPDX-License-Identifier: LPPL-1.3c
% Copyright (C) 2026 Christophe Jorssen
% Author and Current Maintainer: Christophe Jorssen <christophe.jorssen@gmail.com>
% LPPL maintenance status: maintained
% This file is part of luafemm. See LICENSE and LICENSES.md for its terms.
\section{Axisymmetric mapped exterior regions}
\label{sec:exterior}

A Kelvin exterior represents the space outside a physical sphere by a separate
computational sphere. Both appear as half-disks in the $(r,z)$ section. The
mapped sphere is not an additional physical body. Its scale and position serve
the computation; field coordinates and volume averages inside it are mapped
quantities and must not be presented as physical-space measurements.

\begin{key}{/femm/problem/exterior inner radius=\meta{number} (initially 0)}
Physical sphere radius in model units. Positive when a Kelvin region is used.
\end{key}
\begin{key}{/femm/problem/exterior outer radius=\meta{number} (initially 0)}
Mapped sphere radius in model units. Despite the inherited FEMM terminology,
it may be smaller than the physical radius. No ordering of the two radii is required.
\end{key}
\begin{key}{/femm/problem/exterior center z=\meta{number} (initially 0)}
Mapped sphere's axial centre in model units. Both spheres are centred on $r=0$.
\end{key}
\begin{key}{/femm/region/exterior=none\textbar kelvin (initially none)}
Mark only the mapped region. It must be linear, isotropic, nonconducting and
source-free, with no permanent magnetization. Ordinary static and harmonic
axisymmetric models are supported, not transient or operating-point responses.
At an element centroid the relative reluctivity is multiplied by
$(r^2+(z-z_o)^2)R_i/R_o^3$, following FEMM's exterior coefficient convention.
\end{key}
\begin{key}{/femm/pair/scale=\meta{positive number} (initially 1)}
For a Kelvin interface, scale coordinates from the first frontier to the second.
Use zero rotation and radial shift, with the required axial shift. The constraint
preserves $\psi=2\pi rA_\phi$: if $r_2=s r_1$, then $A_{\phi,1}=s A_{\phi,2}$
for periodic pairing. Ordinary planar pairs still require scale 1.
\end{key}

Declare the compound physical domain, currents/materials and mapped region,
then the two |type=periodic| frontiers and their pair. Use exactly matching
evaluated traces at the declared scale. The example below uses shared polygon
definitions so all domain, region and frontier captures agree. Refining the
angular polygon and the field mesh are independent operations.

% The complete generic scene is included as executed source below.

\begin{codeexample}[width=0pt]
% Generic body: physical hemisphere and separate Kelvin computational bubble.
\begingroup
\def\Physical{(0,-20) \foreach \a in {-84,-78,...,84}{--({20*cos(\a)},{20*sin(\a)})} --(0,20)}
\def\Mapped{(0,30) \foreach \a in {-84,-78,...,84}{--({5*cos(\a)},{35+5*sin(\a)})} --(0,40)}
\tikzpicture[femm/problem={type=axisymmetric,xmin=0,xmax=20,
 ymin=-20,ymax=40,mesh size=2,exterior inner radius=20,
 exterior outer radius=5,exterior center z=35},scale=1.4]
 \path[femm/domain] \Physical--cycle
   \Mapped--cycle;
 \path[femm/region={material=air,current density=1e6}] (5,-3) rectangle (10,3);
 \path[femm/region={material=air,exterior=kelvin}]
   \Mapped--cycle;
 \path[femm/frontier={name=physical,type=periodic,curve tolerance=.03}]
   \Physical;
 \path[femm/frontier={name=mapped,type=periodic,curve tolerance=.0075}]
   \Mapped;
 \femmpair[first=physical,second=mapped,scale=.25,shift y=35,tolerance=.03]
 \draw[gray,dashed] \Physical--cycle \Mapped--cycle;
 \pic[femm/lines=10] {femm field};
 \draw[orange,thick] (5,-3) rectangle (10,3);
 \node[right] at (5,35) {Mapped exterior};
 \draw[->] (0,-22)--(0,42) node[above] {$z$};
\endtikzpicture
\par
\endgroup
\end{codeexample}

Native FEMM imports retain exterior radii, centre and label flags. Equal-angle
periodic arcs use synchronized angular subdivisions. Static or harmonic ANS
answers retain their nodal flux constraints; the mapped answer can therefore
be restored without solving. Exporting scaled interfaces as an editable FEM
would require unequal-length periodic straight primitives, so that export is
rejected. ANS export stores the constraints in its nodal-pair trailer; extracting
its embedded polygonal geometry alone is not an equivalent editable problem.
Windows GUI acceptance remains unverified.

Real point-potential properties on retained geometric vertices are also imported
and exported. Point currents, complex prescribed point potentials and nonzero
axisymmetric point-flux properties are explicitly rejected. An isolated input
point which the mesher did not retain cannot silently become a nearby constraint.
